\documentclass[12pt,a4paper]{article}

\usepackage{geometry}
\geometry{a4paper, margin=2.5cm}

\usepackage{graphicx}
\usepackage{xeCJK}
\setCJKmainfont{PingFang TC} % Standard macOS Traditional Chinese font

\usepackage{setspace}
\usepackage{titlesec}
\usepackage{amssymb}

\setstretch{1.6}

\titleformat{\section}{\fontsize{18}{20}\bfseries}{\thesection}{1em}{}
\titleformat{\subsection}{\fontsize{16}{18}\bfseries}{\thesubsection}{1em}{}

\begin{document}

\begin{center}
{\LARGE 青年百億海外圓夢基金計畫}\\
\vspace{0.5cm}
{\Large 海外翱翔組 計畫申請書}\\
\vspace{1cm}

申請人：陳以珊\\
申請計畫：AI 演算法遠征\\
申請年度：115 年度第二梯次
\end{center}

\newpage

\section{申請青年之學經歷}

本人目前就讀於上海交通大學電子與計算機工程學系（Electrical and Computer Engineering），預計於2026年畢業。在學期間，我的研究與學習方向主要集中於人工智慧、機器學習以及資料科學的應用。

在課程方面，我修習了 Computational Methods 以及 Data Science and Analytics using Python 等資料科學相關課程，並透過多項專案實作深入理解機器學習模型的建構流程與應用場景。

在學術發展上，我曾前往日本慶應義塾大學進行交換學習，並在此前於法國特魯瓦科技大學以及關西學院大學進行短期交流。這些跨文化學習經驗使我能夠在多元的學術環境中培養國際視野，並提升跨文化溝通能力。

\subsection{人工智慧與資料科學專案經驗}

在過去的專案經驗中，我主要參與並主導多項與人工智慧相關的研究與開發工作。

其中一項重要專案為基於 Graph Neural Network 與 Transformer 的材料性質預測系統。該專案基於 Open Catalyst 2022（OC22）資料集，透過圖神經網路建模材料的原子結構，並結合 Transformer 架構以捕捉長距離結構關係。

\begin{figure}[h]
\centering
\includegraphics[width=0.8\textwidth]{model_architecture.png}
\caption{GNN + Transformer 材料性質預測模型架構（示意圖）}
\end{figure}

透過這一研究，我建立了完整的機器學習 pipeline，包括資料清理、特徵工程、模型訓練與性能評估。最終模型在驗證資料集上達到 RMSE 0.38 的預測表現。

\subsection{Agentic AI 系統開發}

除了機器學習模型研究外，我亦致力於 AI 系統架構的設計與實現。目前我正在開發一個名為 terny.ai 的科研機會匹配系統。

該系統透過自動化爬蟲與大型語言模型解析全球實驗室與研究機會資訊，並透過多智能體（Multi-Agent）架構進行學生能力建模與科研方向匹配。

\begin{figure}[h]
\centering
\includegraphics[width=0.8\textwidth]{agent_architecture.png}
\caption{Agentic AI 系統架構示意圖}
\end{figure}

此系統整合了資料擷取、語意解析、任務規劃與決策模組，形成完整的 AI agent workflow。

\section{對所申請計畫的理解}

Lambton College 的 Digital Transformation Lab 在人工智慧研究與產業應用方面具有顯著成果。近年來，該實驗室積極推動 Responsible AI 的研究方向，並推出 Ethical and Responsible AI Tool。

該工具基於加拿大政府提出的 FASTER 原則，強調人工智慧系統的公平性、透明性以及可問責性。

\begin{figure}[h]
\centering
\includegraphics[width=0.8\textwidth]{faster_framework.png}
\caption{FASTER Responsible AI Framework}
\end{figure}

透過這一框架，研究人員能夠在 AI 開發的早期階段即評估資料與模型中的潛在偏差，並確保 AI 系統符合倫理與治理標準。

\section{申請動機}

在過去的研究與專案經驗中，我逐漸意識到人工智慧技術不僅是一種計算方法，更是一種能夠深刻影響社會與產業結構的重要工具。

然而，AI 系統在實際應用中也面臨多項挑戰，例如資料偏見、模型可解釋性以及系統責任歸屬等問題。因此，我對於 Responsible AI 的研究方向產生了濃厚興趣。

Lambton College 的 Digital Transformation Lab 正是在這一領域具有領先研究的機構之一。透過參與 AI 演算法遠征計畫，我希望能夠深入了解 AI 系統從研究到產業落地的完整流程。

\section{計畫執行規劃}

本計畫的學習與研究將分為三個階段進行。

第一階段為研究環境與技術架構的熟悉，包括了解研究團隊的 AI 系統架構、資料處理流程以及既有模型。

第二階段為模型開發與實驗設計。我將參與機器學習模型的開發與優化工作，包括特徵工程、模型訓練與性能評估。

第三階段為成果整理與系統優化。透過實驗結果分析，我將進一步探索 AI 模型在實際應用中的性能與限制。

\section{回國後之規劃}

完成本計畫後，我希望能將在海外學習到的 AI 技術與研究經驗帶回台灣。

首先，我計畫在校園內舉辦 AI 技術分享活動，向學生介紹 AI 系統開發與 Responsible AI 的概念。

其次，我希望將 Responsible AI 的理念融入未來的研究與系統設計中，推動更具倫理與社會責任的 AI 技術發展。

最後，我也希望透過建立跨國研究合作網絡，促進台灣與國際研究機構之間的交流與合作。

\section{經費需求}

本計畫之經費需求依據官方公告之合作計畫編列。

\begin{itemize}
\item [$\boxtimes$] 同意上述辦理方式
\item [$\square$] 不同意上述辦理方式
\end{itemize}

\section{外語能力}

本人曾於日本與法國進行交換學習，並具備良好的英語溝通能力，能夠在國際研究環境中進行學術交流與合作。

\section{其他證明文件}

附件包括：

\begin{itemize}
\item 履歷
\item 專案成果
\item 獎學金證明
\item 推薦信
\end{itemize}

\end{document}